% Guion de la práctica: ensayos de la máquina asíncrona trifásica (VEMLAB)
% Compilar con:  tectonic guion_maquina_asincrona.tex
\documentclass[11pt,a4paper]{article}

\usepackage{fontspec}
\usepackage[spanish,es-noshorthands,es-tabla]{babel}
\usepackage[a4paper,margin=2.2cm,headheight=14pt]{geometry}
\usepackage{amsmath,amssymb}
\usepackage{booktabs,tabularx,array,multirow}
\usepackage{enumitem}
\usepackage{xcolor}
\usepackage{tikz}
\usetikzlibrary{calc,positioning}
\usepackage{fancyhdr}
\usepackage{titlesec}
\usepackage{tcolorbox}
\usepackage{hyperref}


\definecolor{acento}{HTML}{1F5FA8}
\definecolor{aviso}{HTML}{B23A1E}
\hypersetup{colorlinks=true,urlcolor=acento,linkcolor=acento}

\titleformat{\section}{\sffamily\large\bfseries\color{acento}}{\thesection.}{0.6em}{}
\titleformat{\subsection}{\sffamily\normalsize\bfseries}{\thesubsection.}{0.5em}{}
\titlespacing*{\section}{0pt}{14pt}{6pt}

\pagestyle{fancy}
\fancyhf{}
\lhead{\sffamily\small Máquinas Eléctricas · Laboratorio}
\rhead{\sffamily\small Práctica: máquina asíncrona trifásica}
\cfoot{\sffamily\small\thepage}
\renewcommand{\headrulewidth}{0.4pt}

\setlist{itemsep=2pt,topsep=4pt}
\setlength{\parindent}{0pt}
\setlength{\parskip}{5pt}

\newtcolorbox{seguridad}{colback=aviso!6,colframe=aviso,boxrule=0.6pt,arc=2pt,
  left=6pt,right=6pt,top=4pt,bottom=4pt,fonttitle=\sffamily\bfseries,title=Seguridad}
\newtcolorbox{nota}{colback=acento!5,colframe=acento!70,boxrule=0.5pt,arc=2pt,
  left=6pt,right=6pt,top=4pt,bottom=4pt}

\renewcommand{\arraystretch}{1.35}
\newcolumntype{C}{>{\centering\arraybackslash}X}

\begin{document}

% ---------------------------------------------------------------- portada
\begin{center}
  {\sffamily\small\color{acento} LABORATORIO VIRTUAL VEMLAB}\\[4pt]
  {\sffamily\LARGE\bfseries Ensayos de la máquina asíncrona trifásica}\\[4pt]
  {\sffamily\large Vacío, cortocircuito y carga · Circuito equivalente}
\end{center}

\vspace{6pt}
\begin{tabularx}{\textwidth}{|l|X|l|X|}
  \hline
  \textsf{Alumno/a} & & \textsf{Grupo} & \\ \hline
  \textsf{Fecha} & & \textsf{Máquina ensayada} & \\ \hline
\end{tabularx}

% ---------------------------------------------------------------- 1
\section{Objetivos}
\begin{itemize}
  \item Interpretar la placa de características y conectar correctamente la caja de bornes a la red disponible.
  \item Medir la potencia trifásica por el método de los dos vatímetros (montaje Aron).
  \item Realizar los ensayos de vacío y de cortocircuito (rotor bloqueado) y obtener a partir de ellos los parámetros del circuito equivalente por fase.
  \item Realizar el ensayo en carga y obtener las curvas de funcionamiento: corriente, factor de potencia, rendimiento y deslizamiento frente a la potencia útil.
  \item Contrastar los resultados medidos con los que predice el circuito equivalente.
\end{itemize}

% ---------------------------------------------------------------- 2
\section{Material del puesto de trabajo}
El banco reproduce un puesto real de laboratorio. Abre VEMLAB en el navegador y elige en la barra superior la máquina que te indique el profesor.

\begin{tabularx}{\textwidth}{@{}>{\bfseries}l X@{}}
  \toprule
  Fuente trifásica regulable & Autotransformador 0–440\,V entre fases, 50\,Hz, 16\,A. Interruptor automático Q1 (C16), pilotos de red y disparo y seta de parada de emergencia. Bornes L1, L2, L3 y N. \\
  Vatímetros W1 y W2 & Digitales, 20\,A / 600\,V. Circuito amperimétrico (I*, I) y voltimétrico (U*, U); el asterisco indica el borne de polaridad. Indican el signo de la potencia. \\
  Amperímetro & 20\,A c.a., bornes A y COM. \\
  Voltímetro & 600\,V c.a., bornes V y COM. \\
  Polímetro & En posición óhmetro (Ω). Solo se usa \emph{sin tensión}. \\
  Caja de bornes & U1, V1, W1 / W2, U2, V2 con chapas de puente para conexión estrella o triángulo. \\
  Freno & Consigna de par regulable, bloqueo del rotor, lectura de par (N·m) y velocidad (min$^{-1}$), indicador de sentido de giro y sondas térmicas (PTC) de los devanados. \\
  Osciloscopio / registrador & Tendencia de las medidas, ondas $u$–$i$ instantáneas y curva par–velocidad (X–Y). \\
  \bottomrule
\end{tabularx}

\begin{seguridad}
  \begin{itemize}[leftmargin=12pt]
    \item Todo el montaje (cables y puentes) se hace y se modifica \textbf{con Q1 abierto}. El banco no permite tocar el cableado con tensión.
    \item Antes de cerrar Q1, lleva el autotransformador a \textbf{0\,V} y sube la tensión poco a poco.
    \item En el ensayo de cortocircuito no superes la corriente nominal. A tensión nominal con el rotor bloqueado la corriente sería de 5 a 7 veces la nominal.
    \item Ante un cortocircuito, Q1 dispara. Si los devanados se calientan demasiado, las sondas PTC desconectan el motor. Anota en el informe cualquier disparo y su causa.
  \end{itemize}
\end{seguridad}

% ---------------------------------------------------------------- 3
\section{Placa de características y conexión}
Copia los datos de la placa del motor (está sobre la carcasa) y calcula las magnitudes derivadas.

\begin{tabularx}{\textwidth}{|X|X|X|X|X|X|}
  \hline
  $P_n$ (kW) & $U_n$ Δ/Y (V) & $I_n$ Δ/Y (A) & $n_n$ (min$^{-1}$) & $\cos\varphi_n$ & $f$ (Hz) \\ \hline
  & & & & & \\[6pt] \hline
  \multicolumn{2}{|l|}{Pares de polos $p$ =} & \multicolumn{2}{l|}{$n_1 = 60f/p$ =} & \multicolumn{2}{l|}{$s_n = (n_1-n_n)/n_1$ =} \\[4pt] \hline
  \multicolumn{3}{|l|}{Par nominal $T_n = P_n/\omega_n$ =} & \multicolumn{3}{l|}{Rendimiento nominal $\eta_n$ =} \\[4pt] \hline
\end{tabularx}

La red del laboratorio es de \textbf{400\,V entre fases}. Cada bobina debe quedar a su tensión nominal, que es la menor de las dos tensiones de la placa.

\textbf{Conexión elegida:} \quad $\square$ Estrella (Y) \qquad $\square$ Triángulo (Δ) \qquad Justificación: \hrulefill

\vspace{2pt}
\begin{nota}
  Con estrella, la tensión de cada bobina es $U/\sqrt3$. Con triángulo, es $U$. Un motor 230/400\,V conectado en triángulo a 400\,V soportaría $\sqrt3$ veces su tensión nominal.
\end{nota}

% ---------------------------------------------------------------- 4
\section{Medida de la resistencia del estátor}
Con Q1 abierto, \textbf{quita los puentes} de la caja de bornes y mide con el polímetro la resistencia de cada bobina.

\begin{tabularx}{\textwidth}{|X|X|X|X|X|}
  \hline
  $R_{U1\text{–}U2}$ (Ω) & $R_{V1\text{–}V2}$ (Ω) & $R_{W1\text{–}W2}$ (Ω) & $R_\text{bobina}$ media (Ω) & $R_1$ estrella equiv. (Ω) \\ \hline
  & & & & \\[6pt] \hline
\end{tabularx}

Todos los cálculos de este guion se hacen \textbf{por fase de la estrella equivalente}. Si el motor trabaja en estrella, $R_1 = R_\text{bobina}$. Si trabaja en triángulo, $R_1 = R_\text{bobina}/3$. En un laboratorio real habría que corregir esta resistencia a la temperatura de funcionamiento. En el simulador no hace falta.

Vuelve a colocar los puentes según la conexión elegida en el apartado anterior.

% ---------------------------------------------------------------- 5
\section{Montaje: método de los dos vatímetros (Aron)}
El mismo montaje sirve para los tres ensayos.
\begin{itemize}
  \item Circuito amperimétrico de W1 en serie con L1, y de W2 en serie con L2. La corriente entra por I*.
  \item Circuito voltimétrico de W1 entre L1 (U*) y L3 (U). El de W2, entre L2 (U*) y L3 (U).
  \item El amperímetro va en serie con una fase y el voltímetro entre dos fases, en bornes del motor.
\end{itemize}

\begin{center}
\begin{tikzpicture}[line width=0.7pt, font=\sffamily\small,
    inst/.style={draw, circle, minimum size=7mm, inner sep=0pt, fill=white},
    bob/.style={draw, rectangle, minimum width=11mm, minimum height=4mm, fill=white},
    vbob/.style={draw, rectangle, minimum width=4mm, minimum height=6.5mm, fill=white}]
  % líneas
  \foreach \n/\y in {L1/2.4, L2/1.2, L3/0} {
    \node[left] at (0,\y) {\n};
    \draw (0,\y) -- (10.2,\y);
    \fill (0,\y) circle (1.6pt);
  }
  % amperímetro en L1
  \node[inst] at (1.3,2.4) {A};
  % W1
  \node[bob] (w1i) at (3.2,2.4) {};
  \node[above=3pt] at (w1i.north) {W1};
  \node at ($(w1i.west)+(-0.12,0.22)$) {*};
  % voltimétrica de W1: L1 -> L3 (cruza L2 sin conexión)
  \draw (4.3,2.4) -- (4.3,0);
  \node[vbob] at (4.3,1.8) {};
  \fill (4.3,2.4) circle (1.6pt); \fill (4.3,0) circle (1.6pt);
  \node at (4.58,2.2) {*};
  \draw[densely dashed, gray] (w1i.east) -- (4.3,2.12);
  % W2
  \node[bob] (w2i) at (5.6,1.2) {};
  \node[above=3pt] at (w2i.north) {W2};
  \node at ($(w2i.west)+(-0.12,0.22)$) {*};
  \draw (6.7,1.2) -- (6.7,0);
  \node[vbob] at (6.7,0.6) {};
  \fill (6.7,1.2) circle (1.6pt); \fill (6.7,0) circle (1.6pt);
  \node at (6.98,1.0) {*};
  \draw[densely dashed, gray] (w2i.east) -- (6.7,0.92);
  % voltímetro entre L1 y L2
  \draw (8.2,2.4) -- (8.2,2.15); \draw (8.2,1.45) -- (8.2,1.2);
  \node[inst] at (8.2,1.8) {V};
  \fill (8.2,2.4) circle (1.6pt); \fill (8.2,1.2) circle (1.6pt);
  % motor
  \draw (10.2,2.4) -- (10.85,2.4) -- (10.85,1.8);
  \draw (10.2,1.2) -- (10.6,1.2);
  \draw (10.2,0) -- (10.85,0) -- (10.85,0.6);
  \node[draw, circle, minimum size=17mm, fill=white, align=center] at (11.45,1.2) {M\\[-1pt]3\,$\sim$};
  \node[font=\sffamily\scriptsize] at (10.25,2.6) {U1};
  \node[font=\sffamily\scriptsize] at (10.25,1.4) {V1};
  \node[font=\sffamily\scriptsize] at (10.25,0.2) {W1};
\end{tikzpicture}
\end{center}

\vspace{-4pt}
Con este montaje, la potencia absorbida y el factor de potencia son:
\[
  P = W_1 + W_2, \qquad \tan\varphi = \sqrt3\,\frac{W_1 - W_2}{W_1 + W_2}, \qquad \cos\varphi = \frac{P}{\sqrt3\,U I}.
\]
Si $\cos\varphi < 0{,}5$, uno de los vatímetros marca una potencia \textbf{negativa}. Pasa siempre en vacío. Se suma con su signo.

\begin{nota}
  Antes de empezar, comprueba el montaje con el motor en marcha. La suma $W_1 + W_2$ debe ser coherente, el amperímetro debe marcar lo mismo en cualquier fase y el voltímetro, unos 400\,V a tensión nominal.
\end{nota}

% ---------------------------------------------------------------- 6
\section{Ensayo de vacío}
\subsection*{Procedimiento}
\begin{enumerate}
  \item Rotor libre y consigna del freno a 0\,N·m. Autotransformador a 0\,V. Cierra Q1.
  \item Sube la tensión hasta la nominal. El motor arranca y gira casi a la velocidad de sincronismo $n_1$.
  \item Sube hasta $\approx 1{,}1\,U_n$ y toma medidas \textbf{bajando} la tensión por escalones. Detente cuando la velocidad empiece a caer de forma apreciable.
\end{enumerate}

\begin{tabularx}{\textwidth}{|C|C|C|C|C|C|C|C|}
  \hline
  $U_0$ (V) & $I_0$ (A) & $W_1$ (W) & $W_2$ (W) & $P_0$ (W) & $\cos\varphi_0$ & $n$ (min$^{-1}$) & $P_0 - 3I_0^2R_1$ (W) \\ \hline
   & & & & & & & \\[5pt] \hline
   & & & & & & & \\[5pt] \hline
   & & & & & & & \\[5pt] \hline
   & & & & & & & \\[5pt] \hline
   & & & & & & & \\[5pt] \hline
   & & & & & & & \\[5pt] \hline
   & & & & & & & \\[5pt] \hline
   & & & & & & & \\[5pt] \hline
   & & & & & & & \\[5pt] \hline
\end{tabularx}

\subsection*{Cálculos}
\begin{enumerate}
  \item Representa $P_{Fe} + P_m = P_0 - 3 I_0^2 R_1$ frente a $U_0^2$. Ajusta una recta y extrapola a $U_0 = 0$. La ordenada en el origen son las pérdidas mecánicas $P_m$.
  \item Con la medida a tensión nominal, calcula las pérdidas en el hierro $P_{Fe} = P_0 - 3I_0^2R_1 - P_m$.
  \item Calcula los parámetros de la rama paralelo por fase. Usa solo el punto de tensión nominal, porque la máquina se satura:
  \[
    R_{Fe} = \frac{(U_n/\sqrt3)^2}{P_{Fe}/3} = \frac{U_n^2}{P_{Fe}}, \qquad
    X_\mu \approx \frac{U_n/\sqrt3}{I_0 \sin\varphi_0}.
  \]
  \item Representa también $I_0$ y $\cos\varphi_0$ frente a $U_0$ y comenta su forma.
\end{enumerate}

% ---------------------------------------------------------------- 7
\section{Ensayo de cortocircuito (rotor bloqueado)}
\begin{seguridad}
  Activa el \textbf{bloqueo del rotor} con el motor parado. Comprueba que el autotransformador está a 0\,V \emph{antes} de cerrar Q1. Sube la tensión despacio, con los botones $-$/$+$, y vigila el amperímetro. Haz el ensayo rápido y observa el indicador de temperatura de los devanados.
\end{seguridad}

\subsection*{Procedimiento}
\begin{enumerate}
  \item Mantén el mismo montaje y activa el bloqueo del rotor.
  \item Cierra Q1 y sube la tensión hasta que la corriente sea aproximadamente la nominal, $I_{cc} \approx I_n$. Toma varios puntos por el camino.
  \item El freno mide el par de arranque correspondiente a cada tensión.
\end{enumerate}

\begin{tabularx}{\textwidth}{|C|C|C|C|C|C|C|}
  \hline
  $U_{cc}$ (V) & $I_{cc}$ (A) & $W_1$ (W) & $W_2$ (W) & $P_{cc}$ (W) & $\cos\varphi_{cc}$ & $T_\text{arr}$ (N·m) \\ \hline
   & & & & & & \\[5pt] \hline
   & & & & & & \\[5pt] \hline
   & & & & & & \\[5pt] \hline
   & & & & & & \\[5pt] \hline
   & & & & & & \\[5pt] \hline
   & & & & & & \\[5pt] \hline
\end{tabularx}

\subsection*{Cálculos (punto con $I_{cc} = I_n$)}
\[
  R_{cc} = \frac{P_{cc}}{3 I_{cc}^2} = R_1 + R'_2, \qquad
  Z_{cc} = \frac{U_{cc}/\sqrt3}{I_{cc}}, \qquad
  X_{cc} = \sqrt{Z_{cc}^2 - R_{cc}^2} = X_1 + X'_2 .
\]
De aquí se obtiene $R'_2 = R_{cc} - R_1$. Si no hay más información, se toma $X_1 = X'_2 = X_{cc}/2$.
\begin{enumerate}
  \item Representa $U_{cc}$ frente a $I_{cc}$ y $P_{cc}$ frente a $I_{cc}^2$. ¿Son rectas? ¿Por qué?
  \item Estima la corriente y el par de arranque a tensión nominal: $I_\text{arr} \approx I_{cc}\,U_n/U_{cc}$ y $T_\text{arr} \propto U^2$. Compáralos con $I_n$ y $T_n$.
\end{enumerate}

% ---------------------------------------------------------------- 8
\section{Ensayo en carga}
\subsection*{Procedimiento}
\begin{enumerate}
  \item Desactiva el bloqueo del rotor y pon la consigna del freno a 0\,N·m. Arranca el motor y ajusta $U = 400$\,V.
  \item Aumenta el par del freno por escalones, desde vacío hasta $\approx 1{,}25\,T_n$. En cada escalón, reajusta la tensión a 400\,V si hace falta y anota las medidas.
\end{enumerate}

\begin{tabularx}{\textwidth}{|C|C|C|C|C|C|C||C|C|C|C|}
  \hline
  $U$ (V) & $I$ (A) & $W_1$ (W) & $W_2$ (W) & $P_1$ (W) & $n$ (min$^{-1}$) & $T$ (N·m) & $s$ (\%) & $P_u$ (W) & $\eta$ (\%) & $\cos\varphi$ \\ \hline
   & & & & & & & & & & \\[5pt] \hline
   & & & & & & & & & & \\[5pt] \hline
   & & & & & & & & & & \\[5pt] \hline
   & & & & & & & & & & \\[5pt] \hline
   & & & & & & & & & & \\[5pt] \hline
   & & & & & & & & & & \\[5pt] \hline
   & & & & & & & & & & \\[5pt] \hline
   & & & & & & & & & & \\[5pt] \hline
   & & & & & & & & & & \\[5pt] \hline
\end{tabularx}

\subsection*{Cálculos}
\[
  P_1 = W_1 + W_2,\quad P_u = T\,\omega = T\,\frac{\pi n}{30},\quad \eta = \frac{P_u}{P_1},\quad s = \frac{n_1 - n}{n_1},\quad \cos\varphi = \frac{P_1}{\sqrt3\,U I}.
\]
\begin{enumerate}
  \item Representa $I$, $\eta$, $\cos\varphi$ y $s$ frente a $P_u$. Señala en cada gráfica el punto nominal de la placa.
  \item Representa $T$ frente a $n$ en la zona de funcionamiento estable.
  \item Haz el balance de potencias en el punto más próximo al nominal: $P_1 = P_{cu1} + P_{Fe} + P_{cu2} + P_m + P_u$, con $P_{cu2} = s\,P_a$ y $P_a = P_1 - P_{cu1} - P_{Fe}$.
\end{enumerate}

% ---------------------------------------------------------------- 9
\section{Curva par–velocidad con el osciloscopio}
El registrador puede trazar la curva completa $T$–$n$ por el \emph{método de aceleración}. Durante el arranque libre, el par motor es $T = J\,\mathrm{d}\omega/\mathrm{d}t + T_\text{rozamiento}$.
\begin{enumerate}
  \item Selecciona el modo \textbf{Par–velocidad} en el osciloscopio. Pon la consigna del freno a 0 y el motor parado.
  \item Arranca el motor, preferiblemente a tensión reducida para que el arranque dure más y la curva tenga más puntos.
  \item Mide sobre la curva el par de arranque, el par máximo y la velocidad a la que se produce. Si has arrancado a tensión reducida, refiérelos a tensión nominal ($T \propto U^2$).
\end{enumerate}

\begin{tabularx}{\textwidth}{|X|X|X|X|}
  \hline
  Tensión de arranque (V) & $T_\text{arr}$ (N·m) & $T_\text{máx}$ (N·m) & $n$ a $T_\text{máx}$ (min$^{-1}$) \\ \hline
  & & & \\[6pt] \hline
\end{tabularx}

% ---------------------------------------------------------------- 10
\section{Resultados: circuito equivalente por fase}
\begin{tabularx}{\textwidth}{|C|C|C|C|C|C|C|}
  \hline
  $R_1$ (Ω) & $R'_2$ (Ω) & $X_1$ (Ω) & $X'_2$ (Ω) & $R_{Fe}$ (Ω) & $X_\mu$ (Ω) & $P_m$ (W) \\ \hline
  & & & & & & \\[8pt] \hline
\end{tabularx}

Dibuja el circuito equivalente con los valores obtenidos.

% ---------------------------------------------------------------- 11
\section{Cuestiones}
\begin{enumerate}
  \item Con el circuito equivalente, calcula para el deslizamiento nominal la corriente, el factor de potencia, el par y el rendimiento. Compáralos con la placa y con el ensayo en carga. Explica las discrepancias.
  \item Calcula con el circuito equivalente el deslizamiento de par máximo $s_m \approx R'_2/\sqrt{R_1^2 + (X_1+X'_2)^2}$ y el par máximo. Compáralos con la curva del osciloscopio.
  \item ¿Por qué uno de los vatímetros marca potencia negativa en vacío? ¿A partir de qué carga pasan los dos a ser positivos?
  \item ¿Qué ocurriría si conectaras el motor a 400\,V con la conexión contraria a la que has elegido? Si quieres, compruébalo en el banco y anota lo que observas.
  \item ¿Cómo se invierte el sentido de giro? Compruébalo con el indicador de sentido del freno.
  \item ¿Por qué el ensayo de vacío solo permite obtener $R_{Fe}$ y $X_\mu$ a la tensión nominal? ¿Qué efecto tiene la saturación sobre $I_0$?
\end{enumerate}

\vfill
\begin{nota}
  \textbf{Entrega.} Informe con las tablas completas, las gráficas pedidas, los cálculos justificados y las respuestas a las cuestiones.
\end{nota}

\end{document}
